\lstset{showstringspaces=false}
\chapter{Dettagli Implementativi e Codici Sorgente}
La presente appendice raccoglie i codici sorgente completi sviluppati in linguaggio Python per la pre-elaborazione dell'immagine WSI. Gli script costituiscono la pipeline automatizzata progettata per trasformare la Whole Slide Image originale in un dataset standardizzato, senza  rumore e pronto per l'addestramento della rete neurale.

\section{Estrazione dei Metadati e Generazione Thumbnail}
Lo script \texttt{task2\_generateJson.py} analizza la WSI tramite la libreria \texttt{openslide}. Il codice estrae le informazioni strutturali fondamentali (livelli di zoom, dimensioni, magnificazione), salvandole in un file JSON, e produce una miniatura a bassa risoluzione (thumbnail) indispensabile per le successive elaborazioni.

\begin{lstlisting}[language=Python, caption={Script completo per l'estrazione dei metadati WSI.}]
	import openslide
	import json
	import glob
	import os
	import sys
	
	# controllo per capire se è stato inserito un nome che esiste di un file della directory
	if len(sys.argv) < 2:
	print("Errore: Devi dirmi quale file cercare!")
	print("Esempio: python3 task2_finale.py CMU... ad esempio")
	exit()
	# se c'è lo salva
	input_utente = sys.argv[1]
	
	files = glob.glob(input_utente)
	
	if not files:
	print(f"ERRORE: Non vedo il file {input_utente} nella cartella!")
	print("Assicurati che il comando 'cp' abbia funzionato.")
	exit()
	
	filename = files[0]
	print(f"file trovato: {filename}")
	print(f"Dimensione: {os.path.getsize(filename) / (1024*1024):.2f} MB")
	
	try:
	# Apre il vetrino
	slide = openslide.OpenSlide(filename)
	
	magnificazione_str = slide.properties.get('openslide.objective-power')
	
	# Prepara i dati per il report
	dati = {
		"nome_file": filename,
		"vendor": slide.properties.get('openslide.vendor'),
		"magnificazione": int(float(magnificazione_str)) if magnificazione_str else "N/A",
		"livelli_zoom": slide.level_count,
		"dimensioni_base": slide.dimensions,
		"nota": "Vetrino caricato con successo in ambiente WSL"
	}
	
	# Salva il JSON con le informazioni presenti in data
	with open('report_wsi.json', 'w') as f:
	json.dump(dati, f, indent=4)
	
	# Crea la Thumbnail
	slide.get_thumbnail((512, 512)).save("thumbnail.png")
	
	print(f"Creato report_wsi.json e thumbnail.png")
	
	
	slide.close()
	
	except Exception as e:
	print(f"Errore durante l'apertura: {e}")
\end{lstlisting}

\section{Segmentazione e creazione della maschera (Metodo di Otsu)}
Il file \texttt{task3\_mask.py} sfrutta la libreria OpenCV per elaborare la miniatura del vetrino. Convertendo l'immagine nello spazio colore HSV rendendola binaria ossia bianco o nero, lo script applica la sogliatura automatica di Otsu sul canale della saturazione, permettendo di creare due macrocategoria, il rumore ossia lo sfondo e il tessuto.

\begin{lstlisting}[language=Python, caption={Script di conversione HSV e applicazione della soglia di Otsu.}]
	import cv2
	import numpy as np
	import matplotlib.pyplot as plt
	import os
	
	image_path = "thumbnail.png"
	
	# Verifica se il file esiste fisicamente prima di caricarlo
	if not os.path.exists(image_path):
	print(f"ERRORE: Il file '{image_path}' non esiste affatto in {os.getcwd()}")
	else:
	img = cv2.imread(image_path)
	if img is None:
	print("ERRORE: Il file esiste ma OpenCV non riesce a leggerlo (forse è corrotto o vuoto).")
	else:
	print("Successo! Immagine caricata correttamente.")
	# Converti l'immagine in HSV (Hue, Saturation, Value)
	# così da poter distinguare il viola del tessuto che ha una Saturazione alta, allo sfondo bianco che ha saturazione zero.
	hsv_image = cv2.cvtColor(img, cv2.COLOR_BGR2HSV)
	# Prendiamo solo il canale Saturazione
	saturazione = hsv_image[:, :, 1] 
	
	# Soglia automatica (Metodo Otsu)
	_, mask = cv2.threshold(saturazione, 0, 255, cv2.THRESH_BINARY + cv2.THRESH_OTSU)
	
	cv2.imwrite("mask_binaria.png", mask)
	
	# 5. creazione grafico delle differenze tra png originale e la maschera binaria
	plt.figure(figsize=(10, 5))
	
	# Immagine originale
	plt.subplot(1, 2, 1)
	plt.imshow(cv2.cvtColor(img, cv2.COLOR_BGR2RGB))
	plt.title("Vetrino Originale (Thumbnail)")
	plt.axis("off")
	
	# Maschera (Bianco = Tessuto, Nero = Sfondo)
	plt.subplot(1, 2, 2)
	plt.imshow(mask, cmap="gray")
	plt.title("Maschera QC (Otsu Threshold)")
	plt.axis("off")
	
	plt.savefig("confronto_qc.png")
	print("creazione dei file mask_binaria.png e confronto_qc.png")
\end{lstlisting}

\section{Estrazione Selettiva delle Patch (Tiling)}
Lo script \texttt{task3\_patching2.py} esegue il mapping proporzionale tra le coordinate della maschera e quelle della slide originale ad alta risoluzione. Attraverso un ciclo iterativo, estrae le patch di dimensione $256 \times 256$ pixel, salvando esclusivamente le porzioni che superano una soglia minima di densità tissutale (75\%).

\begin{lstlisting}[language=Python, caption={Script per il patching proporzionale basato su maschera di densità.}]
	import openslide
	import cv2
	import numpy as np
	import os
	import sys
	import glob
	
	# 1. CONFIGURAZIONE
	PATCH_SIZE = 256        
	LEVEL_TO_READ = 0       
	SOGLIA_DENSITA = 191    
	LIMITATORE = 100        
	
	# 2. CARICAMENTO FILE
	mask_path = 'mask_binaria.png' # Definiamo il percorso mancante
	mask = cv2.imread(mask_path, cv2.IMREAD_GRAYSCALE)
	
	if mask is None:
	print(f"ERRORE: Non esiste {mask_path}!")
	sys.exit()
	
	if len(sys.argv) < 2:
	print("Errore: Passami il nome del vetrino (es. CMU-1.tiff)")
	sys.exit()
	
	wsi_path = glob.glob(sys.argv[1])[0]
	slide = openslide.OpenSlide(wsi_path)
	
	# Calcolo scale
	w_slide, h_slide = slide.dimensions
	h_mask, w_mask = mask.shape
	print(f"Vetrino: {w_slide}x{h_slide} | Maschera: {w_mask}x{h_mask}")
	
	output_dir = "patches"
	os.makedirs(output_dir, exist_ok=True)
	
	# 3. CICLO DI ESTRAZIONE
	print("Inizio estrazione...")
	patches_saved = 0
	
	for y in range(0, h_slide - PATCH_SIZE, PATCH_SIZE):
	for x in range(0, w_slide - PATCH_SIZE, PATCH_SIZE):
	
	# Trasformiamo x, y del vetrino gigante in coordinate per la maschera piccola
	m_x = int(x / w_slide * w_mask)
	m_y = int(y / h_slide * h_mask)
	m_size = int(PATCH_SIZE / w_slide * w_mask)
	
	# ESTRAIAMO LA PATCH DALLA MASCHERA
	# Ora x e y esistono perché siamo dentro il ciclo!
	mask_patch = mask[m_y : m_y + m_size, m_x : m_x + m_size]
	
	# CONTROLLO DENSITÀ
	if mask_patch.size > 0 and np.mean(mask_patch) > SOGLIA_DENSITA:
	
	# Se la patch e' "densa", la leggiamo ad alta risoluzione
	patch = slide.read_region((x, y), LEVEL_TO_READ, (PATCH_SIZE, PATCH_SIZE))
	patch = patch.convert("RGB")
	
	patch_name = f"{output_dir}/patch_{x}_{y}.png"
	patch.save(patch_name)
	patches_saved += 1
	
	if patches_saved >= LIMITATORE:
	break
	if patches_saved >= LIMITATORE:
	break
	
	print(f"\nSalvate {patches_saved} patch dense in '{output_dir}'.")
\end{lstlisting}

\section{Controllo Qualità: Analisi della Nitidezza}
Per garantire input validi alla rete neurale, il codice \texttt{task3\_analysis.py} implementa un controllo qualità basato sull'operatore Laplaciano. Lo script calcola la varianza dei pixel per identificare e scartare le immagini sfocate, salvando un report dettagliato in formato CSV tramite la libreria \texttt{pandas}.

\begin{lstlisting}[language=Python, caption={Script per il filtraggio delle patch sfocate tramite Varianza del Laplaciano.}]
	import cv2
	import numpy as np
	import os
	import glob
	import pandas as pd
	import shutil
	
	# creazione cartelle per smistare le patch buona da quelle col rumore in sottofondo
	INPUT_DIR = "patches" # cartella generale
	OUTPUT_DIR_GOOD = "patches_checked/good" # cartella patch buona
	OUTPUT_DIR_BAD = "patches_checked/blur" #cartella patch da non utilizzare
	SOGLIA_BLUR = 100.0  # Sotto questo valore, l'immagine viene considerata sfocata
	
	# vedo se queste due cartelle esistono se no le creo e creo anche quelle sottostanti
	for cartella in [OUTPUT_DIR_GOOD, OUTPUT_DIR_BAD]:
	if not os.path.exists(cartella):
	os.makedirs(cartella)
	
	# Cerca tutte le immagini png nella cartella patches
	files_patch = glob.glob(f"{INPUT_DIR}/*.png")
	
	if not files_patch:
	print(f"Errore: La cartella '{INPUT_DIR}' e' vuota o non esiste!")
	exit()
	
	results = []
	
	# 2. CICLO DI ANALISI
	for file_path in files_patch:
	filename = os.path.basename(file_path)
	img = cv2.imread(file_path)
	
	if img is None:
	continue
	
	# Blur detection (Varianza del Laplaciano)
	gray = cv2.cvtColor(img, cv2.COLOR_BGR2GRAY)
	blur_score = cv2.Laplacian(gray, cv2.CV_64F).var()
	
	# controllo per capire se è un'immagine buona o no e scelta del percoso della cartella
	status = "GOOD"
	if blur_score < SOGLIA_BLUR:
	status = "BLURRY"
	destinazione = os.path.join(OUTPUT_DIR_BAD, filename)
	else:
	destinazione = os.path.join(OUTPUT_DIR_GOOD, filename)
	
	# vede in quale cartella copiare l'immagine appena analizzata
	shutil.copy(file_path, destinazione)
	
	# Salviamo i dati per il report
	results.append({
		"filename": filename,
		"blur_score": round(blur_score, 2),
		"status": status
	})
	
	# SALVATAGGIO REPORT
	df = pd.DataFrame(results)
	csv_name = "report_quality_control.csv"
	df.to_csv(csv_name, index=False)
	
	print("-" * 50)
	print(f"Report salvato in: '{csv_name}'")
	print(f"Patch nitide salvate in: '{OUTPUT_DIR_GOOD}'")
	print(f"Patch sfocate salvate in: '{OUTPUT_DIR_BAD}'")
	print(f"Statistiche:")
	print(f"Totale analizzate: {len(df)}")
	print(f"Accettate (Good): {len(df[df['status']=='GOOD'])}")
	print(f"Scartate (Blurry): {len(df[df['status']=='BLURRY'])}")
\end{lstlisting}

\section{Analisi Visiva Normalizzazione Cromatica}
Lo script \texttt{task3\_normalization.py} è stato sviluppato per validare visivamente il metodo di Macenko. Utilizzando il pacchetto \texttt{torchstain} con backend \texttt{numpy}, separa matematicamente l'Ematossilina dall'Eosina e genera un'immagine di comparazione tra originale, nuclei e normalizzata tramite \texttt{matplotlib}.

\begin{lstlisting}[language=Python, caption={Script per il calcolo e la visualizzazione della standardizzazione cromatica.}]
	import cv2
	import matplotlib.pyplot as plt
	import glob
	import os
	import sys
	
	# setup libreria in quanto a volte mi da errori
	try:
	from torchstain.normalizers import MacenkoNormalizer
	except ImportError:
	try:
	import torchstain
	MacenkoNormalizer = torchstain.normalizers.MacenkoNormalizer
	except Exception as e:
	print("Errore critico libreria torchstain.")
	sys.exit()
	
	# cerca una patch nella cartella delle patch buone
	INPUT_DIR = "patches_checked/good"
	files = glob.glob(f"{INPUT_DIR}/*.png")
	if not files:
	print(f"Errore: Nessuna patch trovata in '{INPUT_DIR}'")
	exit()
	
	img_path = files[0]
	
	# Leggiamo l'immagine come RGB e non BGR
	original_img = cv2.cvtColor(cv2.imread(img_path), cv2.COLOR_BGR2RGB)
	
	try:
	
	normalizer = MacenkoNormalizer(backend='numpy')
	
	# La funzione restituisce: (ImmagineNorm, MatriceH, MatriceE)
	norm_img, H_channel, E_channel = normalizer.normalize(I=original_img, stains=True)
	
	fig, axes = plt.subplots(1, 3, figsize=(15, 5))
	
	axes[0].imshow(original_img)
	axes[0].set_title("1. Originale")
	axes[0].axis('off')
	
	axes[1].imshow(H_channel, cmap='gray')
	axes[1].set_title("2. Ematossilina (Nuclei)")
	axes[1].axis('off')
	
	axes[2].imshow(norm_img)
	axes[2].set_title("3. Normalizzata (Macenko)")
	axes[2].axis('off')
	
	output_file = "confronto_normalizzazione_lib.png"
	plt.savefig(output_file)
	print(f"\n Grafico salvato in: '{output_file}'")
	
	
	except Exception as e:
	print(f"\nErrore durante il calcolo: {e}")
\end{lstlisting}

\section{Pipeline Finale Automatica}
Il file conclusivo \texttt{task3\_final\_pipeline.py} ingloba le fasi di controllo qualità e normalizzazione in un unico flusso iterativo. Analizza le patch grezze, scarta in automatico quelle non idonee basandosi sul Laplaciano, e converte quelle valide tramite l'algoritmo di Macenko, salvando il dataset finale pronto per l'addestramento della ResNet50.

\begin{lstlisting}[language=Python, caption={Pipeline unificata per filtraggio blur e normalizzazione Macenko a livello di dataset.}]
	import cv2
	import numpy as np
	import os
	import glob
	import sys
	import shutil
	
	# --- CONFIGURAZIONE ---
	INPUT_DIR = "patches"              # Dove sono le immagini grezze
	OUTPUT_DIR = "dataset_ready"       # Dove mettiamo quelle pronte per l'AI
	SOGLIA_BLUR = 100.0                # Soglia nitidezza
	
	# Setup Normalizzatore
	try:
	from torchstain.normalizers import MacenkoNormalizer
	normalizer = MacenkoNormalizer(backend='numpy') 
	except:
	import torchstain
	normalizer = torchstain.normalizers.MacenkoNormalizer(backend='numpy')
	
	# Crea cartella output
	if not os.path.exists(OUTPUT_DIR):
	os.makedirs(OUTPUT_DIR)
	
	# Leggi file
	files = glob.glob(f"{INPUT_DIR}/*.png")
	
	count_saved = 0
	count_blur = 0
	count_error = 0
	
	for i, file_path in enumerate(files):
	filename = os.path.basename(file_path)
	
	# Caricamento
	img = cv2.imread(file_path)
	if img is None: continue
	
	# Blur check 
	gray = cv2.cvtColor(img, cv2.COLOR_BGR2GRAY)
	score = cv2.Laplacian(gray, cv2.CV_64F).var()
	
	if score < SOGLIA_BLUR:
	count_blur += 1
	continue # Salta alla prossima
	
	# 4 normalizzazione Macenko
	try:
	# Converti BGR -> RGB per Macenko
	img_rgb = cv2.cvtColor(img, cv2.COLOR_BGR2RGB)
	
	# Normalizza
	norm_img, _, _ = normalizer.normalize(I=img_rgb, stains=True)
	
	# Converti RGB -> BGR per salvare con OpenCV
	final_img = cv2.cvtColor(norm_img, cv2.COLOR_RGB2BGR)
	
	# Salvataggio
	save_path = os.path.join(OUTPUT_DIR, filename)
	cv2.imwrite(save_path, final_img)
	count_saved += 1
	
	except Exception as e:
	print(f"Errore normalizzazione su {filename}: {e}")
	count_error += 1
	
	# Barra di progresso semplice
	sys.stdout.write(f"\r   Progress: {i+1}/{len(files)} | ✅ Salvate: {count_saved} | ðŸ—‘ï¸  Blur: {count_blur}")
	sys.stdout.flush()
\end{lstlisting}